\subsection{Sequence Diagrams}\label{sec:seq-diagrams}
All endpoints must follow RESTful conventions. Use the following templates; each user story below
applies one template as a concrete sequence diagram.

\subsubsection{Resource Creation Flow (Template)}
Use this logic to design your \texttt{POST} and \texttt{PUT} endpoints to ensure consistent error handling
(e.g., handling duplicate entries or missing resources). See Figures~\ref{fig:seq-register},
\ref{fig:seq-todo-create}, and \ref{fig:seq-subtask-create}.

\subsubsection{Resource Read Flow (Template)}
Use this logic to design your \texttt{GET} endpoints. See Figures~\ref{fig:seq-todo-list}
and \ref{fig:seq-subtask-list}.

\subsubsection{Resource Delete Flow (Template)}
Use this logic to design your \texttt{DELETE} endpoints. See Figures~\ref{fig:seq-todo-delete}
and \ref{fig:seq-subtask-delete}.

\subsubsection{Login/Auth Flow (Template)}
Use this logic to design your \texttt{/auth} endpoints to ensure secure credential validation.
See Figure~\ref{fig:seq-login}.

\newpage
\subsubsection{Epic 1: Account Creation}
\textbf{User story:} As a new user, I can register an account to start tracking my todo tasks.
\newline
\texttt{POST /api/auth/register} --- duplicate username/email returns \texttt{409}; validation errors return \texttt{400}.

\begin{figure}[htbp]
  \centering
  \resizebox{\linewidth}{!}{%
  \begin{sequencediagram}
    \newthread{requester}{Client/Requester}
    \newinst[1.3]{api}{API/Interface Layer}
    \newinst[1.3]{logic}{Business Logic Layer}
    \newinst[1.3]{da}{Data Access Layer}
    \newinst[1.3]{store}{Data Storage}
    \begin{sdblock}[gray!10]{SCENARIO A}{[Resource is Unique]}
      \begin{call}{requester}{[Register Request]}{api}{[2xx Success Response]}
        \begin{call}{api}{[Process Registration]}{logic}{[Success Data/Object]}
          \begin{call}[2]{logic}{[Check Existence]}{da}{}
            \begin{call}{da}{[Query]}{store}{[Result (Not Found)]}
            \end{call}
          \end{call}
          \begin{call}[2]{logic}{[Persist New Resource]}{da}{[Resource Object]}
            \begin{call}{da}{[Insert/Write]}{store}{[New Resource Object]}
            \end{call}
          \end{call}
        \end{call}
      \end{call}
    \end{sdblock}
    \begin{sdblock}[gray!10]{SCENARIO B}{[Conflict/Already Exists]}
      \begin{call}{requester}{[Register Request]}{api}{[4xx Error Response]}
        \begin{call}{api}{[Process Registration]}{logic}{[Conflict/Error Signal]}
          \begin{call}[2]{logic}{[Check Existence]}{da}{}
            \begin{call}{da}{[Query]}{store}{[Result (Found)]}
            \end{call}
          \end{call}
        \end{call}
      \end{call}
    \end{sdblock}
  \end{sequencediagram}%
  }
  \caption{Account creation --- \texttt{POST /api/auth/register}}
  \label{fig:seq-register}
\end{figure}


\newpage
\subsubsection{Epic 2: Authentication}
\textbf{User story:} As a user, I can log in and log out to securely access my todo items.
\newline
\texttt{POST /api/auth/login} --- invalid credentials return \texttt{401}.

\begin{figure}[htbp]
  \centering
  \resizebox{\linewidth}{!}{%
  \begin{sequencediagram}
    \newthread{requester}{Client/Requester}
    \newinst[1.1]{api}{API/Interface Layer}
    \newinst[1.1]{logic}{Business Logic Layer}
    \newinst[1.1]{da}{Data Access Layer}
    \newinst[1.1]{util}{Utility/Service}
    \newinst[1.1]{store}{Data Storage}
    \begin{sdblock}[gray!10]{SCENARIO A}{[Success]}
      \begin{call}{requester}{[Login Request]}{api}{[2xx Success Response]}
        \begin{call}{api}{[Execute Process]}{logic}{[Return Success/Token]}
          \begin{call}[2]{logic}{[Request User Data]}{da}{}
            \begin{call}{da}{[Query/Fetch]}{store}{[Return User Object]}
            \end{call}
          \end{call}
          \begin{call}[1]{logic}{[Validate Credentials/Generate Token]}{util}{[Return Result]}
          \end{call}
        \end{call}
      \end{call}
    \end{sdblock}
    \begin{sdblock}[gray!10]{SCENARIO B}{[Failure - Invalid Credentials]}
      \begin{call}{requester}{[Login Request]}{api}{[4xx Error Response]}
        \begin{call}{api}{[Execute Process]}{logic}{[Return Error/Fail]}
          \begin{call}[2]{logic}{[Request User Data]}{da}{}
            \begin{call}{da}{[Query/Fetch]}{store}{[Return Null/Empty]}
            \end{call}
          \end{call}
        \end{call}
      \end{call}
    \end{sdblock}
  \end{sequencediagram}%
  }
  \caption{Authentication --- \texttt{POST /api/auth/login}}
  \label{fig:seq-login}
\end{figure}

\newpage

\subsubsection{Epic 3: Task Management}
\textbf{User story:} As a user, I can create, edit, and delete todo items to keep track of my work.
\newline
\texttt{GET /api/todos} lists todos for the authenticated user; \texttt{POST /api/todos} and
\texttt{PUT /api/todos/\{id\}} create or edit a todo; \texttt{DELETE /api/todos/\{id\}} removes a
todo and cascades its subtasks. Validation errors return \texttt{400}; missing resources return
\texttt{404}.

\begin{figure}[htbp]
  \centering
  \resizebox{\linewidth}{!}{%
  \begin{sequencediagram}
    \newthread{requester}{Client/Requester}
    \newinst[1.3]{api}{API/Interface Layer}
    \newinst[1.3]{logic}{Business Logic Layer}
    \newinst[1.3]{da}{Data Access Layer}
    \newinst[1.3]{store}{Data Storage}
    \begin{call}{requester}{[Request All Todos]}{api}{[2xx Success Response]}
      \begin{call}{api}{[Execute Request]}{logic}{[Return Data]}
        \begin{call}{logic}{[Ask for Data]}{da}{[Return Data]}
          \begin{call}{da}{[Query Data]}{store}{[Return Data]}
          \end{call}
        \end{call}
      \end{call}
    \end{call}
  \end{sequencediagram}%
  }
  \caption{List todos --- \texttt{GET /api/todos}}
  \label{fig:seq-todo-list}
\end{figure}

\begin{figure}[htbp]
  \centering
  \resizebox{\linewidth}{!}{%
  \begin{sequencediagram}
    \newthread{requester}{Client/Requester}
    \newinst[1.3]{api}{API/Interface Layer}
    \newinst[1.3]{logic}{Business Logic Layer}
    \newinst[1.3]{da}{Data Access Layer}
    \newinst[1.3]{store}{Data Storage}
    \begin{sdblock}[gray!10]{SCENARIO A}{[Success]}
      \begin{call}{requester}{[Create or Edit Todo]}{api}{[2xx Success Response]}
        \begin{call}{api}{[Execute Request]}{logic}{[Return Success/Data]}
          \begin{call}{logic}{[Send Todo Data]}{da}{[Return Data]}
            \begin{call}{da}{[Input New Data]}{store}{[Return Data]}
            \end{call}
          \end{call}
        \end{call}
      \end{call}
    \end{sdblock}
    \begin{sdblock}[gray!10]{SCENARIO B}{[Null Message]}
      \begin{call}{requester}{[Create or Edit Todo]}{api}{[4xx Error Response]}
        \begin{call}{api}{[Execute Request]}{logic}{[Return Error]}
          \begin{call}{logic}{[Send Todo Data]}{da}{[Return Null]}
            \begin{call}{da}{[Send to Repository]}{store}{[Result: Null]}
            \end{call}
          \end{call}
        \end{call}
      \end{call}
    \end{sdblock}
  \end{sequencediagram}%
  }
  \caption{Create or edit todo --- \texttt{POST /api/todos}, \texttt{PUT /api/todos/\{id\}}}
  \label{fig:seq-todo-create}
\end{figure}

\begin{figure}[htbp]
  \centering
  \resizebox{\linewidth}{!}{%
  \begin{sequencediagram}
    \newthread{requester}{Client/Requester}
    \newinst[1.3]{api}{API/Interface Layer}
    \newinst[1.3]{logic}{Business Logic Layer}
    \newinst[1.3]{da}{Data Access Layer}
    \newinst[1.3]{store}{Data Storage}
    \begin{call}{requester}{[Delete Todo]}{api}{[2xx Success Response]}
      \begin{call}{api}{[Execute Request]}{logic}{[Return Success]}
        \begin{call}{logic}{[Send Todo ID]}{da}{[Return Success]}
          \begin{call}{da}{[Delete Todo]}{store}{[Return Success]}
          \end{call}
        \end{call}
      \end{call}
    \end{call}
  \end{sequencediagram}%
  }
  \caption{Delete todo --- \texttt{DELETE /api/todos/\{id\}} (API contract: \texttt{204 No Content})}
  \label{fig:seq-todo-delete}
\end{figure}


\newpage


\subsubsection{Epic 4: Subtask Organization}
\textbf{User story:} As a user, I can create, edit, and delete subtask items to better organize my primary tasks.
\newline
\texttt{GET /api/todos/\{id\}} retrieves the parent todo before subtask operations;
\texttt{POST /api/todos/\{id\}/\{subtask\}} and \texttt{PUT /api/todos/\{id\}/\{subtask\}} create or edit
a subtask; \texttt{DELETE /api/todos/\{id\}/\{subtask\}} removes one subtask. The parent todo must
exist and belong to the authenticated user.

\begin{figure}[htbp]
  \centering
  \resizebox{\linewidth}{!}{%
  \begin{sequencediagram}
    \newthread{requester}{Client/Requester}
    \newinst[1.3]{api}{API/Interface Layer}
    \newinst[1.3]{logic}{Business Logic Layer}
    \newinst[1.3]{da}{Data Access Layer}
    \newinst[1.3]{store}{Data Storage}
    \begin{call}{requester}{[Request All Subtasks]}{api}{[2xx Success Response]}
      \begin{call}{api}{[Execute Request]}{logic}{[Return Data]}
        \begin{call}{logic}{[Ask for Data]}{da}{[Return Data]}
          \begin{call}{da}{[Query Data]}{store}{[Return Data]}
          \end{call}
        \end{call}
      \end{call}
    \end{call}
  \end{sequencediagram}%
  }
  \caption{List subtasks --- \texttt{GET /api/todos/\{id\}} (parent todo context)}
  \label{fig:seq-subtask-list}
\end{figure}

\begin{figure}[htbp]
  \centering
  \resizebox{\linewidth}{!}{%
  \begin{sequencediagram}
    \newthread{requester}{Client/Requester}
    \newinst[1.3]{api}{API/Interface Layer}
    \newinst[1.3]{logic}{Business Logic Layer}
    \newinst[1.3]{da}{Data Access Layer}
    \newinst[1.3]{store}{Data Storage}
    \begin{sdblock}[gray!10]{SCENARIO A}{[Success]}
      \begin{call}{requester}{[Create or Edit Subtask]}{api}{[2xx Success Response]}
        \begin{call}{api}{[Execute Request]}{logic}{[Return Success/Data]}
          \begin{call}{logic}{[Send Subtask Data]}{da}{[Return Data]}
            \begin{call}{da}{[Input New Data]}{store}{[Return Data]}
            \end{call}
          \end{call}
        \end{call}
      \end{call}
    \end{sdblock}
    \begin{sdblock}[gray!10]{SCENARIO B}{[Null Message]}
      \begin{call}{requester}{[Create or Edit Subtask]}{api}{[4xx Error Response]}
        \begin{call}{api}{[Execute Request]}{logic}{[Return Error]}
          \begin{call}{logic}{[Send Subtask Data]}{da}{[Return Null]}
            \begin{call}{da}{[Attempt Data Input]}{store}{[Result: Null]}
            \end{call}
          \end{call}
        \end{call}
      \end{call}
    \end{sdblock}
  \end{sequencediagram}%
  }
  \caption{Create or edit subtask --- \texttt{POST /api/todos/\{id\}/\{subtask\}}, \texttt{PUT /api/todos/\{id\}/\{subtask\}}}
  \label{fig:seq-subtask-create}
\end{figure}

\begin{figure}[htbp]
  \centering
  \resizebox{\linewidth}{!}{%
  \begin{sequencediagram}
    \newthread{requester}{Client/Requester}
    \newinst[1.3]{api}{API/Interface Layer}
    \newinst[1.3]{logic}{Business Logic Layer}
    \newinst[1.3]{da}{Data Access Layer}
    \newinst[1.3]{store}{Data Storage}
    \begin{call}{requester}{[Delete Subtask]}{api}{[2xx Success Response]}
      \begin{call}{api}{[Execute Request]}{logic}{[Return Success]}
        \begin{call}{logic}{[Send Subtask ID]}{da}{[Return Success]}
          \begin{call}{da}{[Delete Subtask]}{store}{[Return Success]}
          \end{call}
        \end{call}
      \end{call}
    \end{call}
  \end{sequencediagram}%
  }
  \caption{Delete subtask --- \texttt{DELETE /api/todos/\{id\}/\{subtask\}} (API contract: \texttt{204 No Content})}
  \label{fig:seq-subtask-delete}
\end{figure}

\newpage
